% GENERATED FILE. Edit canonical lessons, then run npm run generate:books.
% canonical-source-hash: fnv1a64:9083777dacff6765
% canonical-lessons: MR-C378-at, MR-C378-apvad, MR-C378-nirbandh, MR-C378-nikas, MR-C378-purtata-karne

\chapter{Conditions and Exceptions}
\label{ch:mr-conditions-and-exceptions}
\hlchaptermodality{378}

\begin{chapteropening}
\textbf{By the end of this chapter:} \emph{I can talk about conditions, exceptions and criteria.}

Complete the last lesson of chapter 378: I can talk about conditions, exceptions and criteria.
\end{chapteropening}

\section[\texorpdfstring{\mr{अट} (aṭ)}{अट (aṭ)}]{\mr{अट} (aṭ) — a condition (stipulation)}

\label{lesson:MR-C378-at}



\begin{quote}
Before the new one: say the Marathi for inspiration, then the Marathi for to bother.
\end{quote}

\subsection*{You'll want to know: \mr{अट}}
\textbf{\mr{अट}} — \emph{aṭ} — \textquotedblleft{}a condition (stipulation)\textquotedblright{}.

\begin{quote}
\mr{करारात} \mr{एक} \mr{महत्त्वाची} \mr{अट} \mr{होती}.
\end{quote}

\begin{quote}
\emph{karārāt ek mahattvācī aṭ hotī.}
\end{quote}

\begin{quote}
There was an important condition in the contract.
\end{quote}

Feminine. Plural: \mr{अटी}, \emph{aṭī}.

\mr{अटीवर}, on condition: \mr{एका} \mr{अटीवर} \mr{हो} \mr{म्हणणे}, to say yes on one condition.

\begin{cousinweb}
\textbf{\mr{अट}} (\emph{aṭ}) is usually traced to the root \emph{*aṭṭ-}, \textquotedblleft{}obstruct, stop\textquotedblright{}, under which CDIAL lists Marathi \emph{aṭ}, \textquotedblleft{}obstruction\textquotedblright{} (CDIAL 179): a condition holds a thing back.
\end{cousinweb}

\begin{grammarlens}[title={What you've built}]
One more word for saying what happens if.
\end{grammarlens}

\subsection*{Your turn}
\begin{itemize}
\raggedright
  \item \emph{Say it:} \emph{aṭ}
  \item \emph{Say it:} \emph{aṭ}, once more
  \item \emph{Say it:} \emph{karārāt ek mahattvācī aṭ hotī.} — There was an important condition in the contract.
  \item \emph{Recall:} say the Marathi for inspiration, then the Marathi for to bother, then say \emph{aṭ} again
\end{itemize}

\begin{tcolorbox}[breakable,colback=teal!4,colframe=teal!35!black,title={Before you move on}]
What does \mr{अट} mean? (\textquotedblleft{}A condition (stipulation)\textquotedblright{}.) Say it once more.
\end{tcolorbox}
\section[\texorpdfstring{\mr{अपवाद} (apvād)}{अपवाद (apvād)}]{\mr{अपवाद} (apvād) — an exception}

\label{lesson:MR-C378-apvad}



\begin{quote}
Before the new one: say the Marathi for to bother, then the Marathi for a condition.
\end{quote}

\subsection*{You'll want to know: \mr{अपवाद}}
\textbf{\mr{अपवाद}} — \emph{apvād} — \textquotedblleft{}an exception\textquotedblright{}.

\begin{quote}
\mr{हा} \mr{नियम} \mr{सगळ्यांना} \mr{लागू} \mr{आहे}, \mr{कोणीही} \mr{अपवाद} \mr{नाही}.
\end{quote}

\begin{quote}
\emph{hā niyam sagaḷyānnā lāgū āhe, koṇīhī apvād nāhī.}
\end{quote}

\begin{quote}
This rule applies to everyone; nobody is an exception.
\end{quote}

Masculine.

\mr{अपवादाने}, as an exception; X \mr{चा} \mr{अपवाद}, the exception to X.

\begin{cousinweb}
\textbf{\mr{अपवाद}} (\emph{apvād}) is a tatsama, taken from Sanskrit \emph{apavāda}, \textquotedblleft{}evil speaking\textquotedblright{}, which grammarians also used for a special rule that sets a general one aside, an exception (Monier-Williams). The same head gave the inherited Marathi \emph{avāī}, \textquotedblleft{}rumour\textquotedblright{} (CDIAL 456).
\end{cousinweb}

\begin{grammarlens}[title={What you've built}]
One more word for saying what happens if.
\end{grammarlens}

\subsection*{Your turn}
\begin{itemize}
\raggedright
  \item \emph{Say it:} \emph{apvād}
  \item \emph{Say it:} \emph{apvād}, once more
  \item \emph{Say it:} \emph{hā niyam sagaḷyānnā lāgū āhe, koṇīhī apvād nāhī.} — This rule applies to everyone; nobody is an exception.
  \item \emph{Recall:} say the Marathi for to bother, then the Marathi for a condition, then say \emph{apvād} again
\end{itemize}

\begin{tcolorbox}[breakable,colback=teal!4,colframe=teal!35!black,title={Before you move on}]
What does \mr{अपवाद} mean? (\textquotedblleft{}An exception\textquotedblright{}.) Say it once more.
\end{tcolorbox}
\section[\texorpdfstring{\mr{निर्बंध} (nirbandh)}{निर्बंध (nirbandh)}]{\mr{निर्बंध} (nirbandh) — a curb; a restriction}

\label{lesson:MR-C378-nirbandh}



\begin{quote}
Before the new one: say the Marathi for a condition, then the Marathi for an exception.
\end{quote}

\subsection*{You'll want to know: \mr{निर्बंध}}
\textbf{\mr{निर्बंध}} — \emph{nirbandh} — \textquotedblleft{}a curb; a restriction\textquotedblright{}.

\begin{quote}
\mr{पुरामुळे} \mr{प्रवासावर} \mr{निर्बंध} \mr{घातले} \mr{आहेत}.
\end{quote}

\begin{quote}
\emph{purāmuḷe pravāsāvar nirbandh ghātle āhet.}
\end{quote}

\begin{quote}
Restrictions have been placed on travel because of the flood.
\end{quote}

Masculine.

Mostly plural: \mr{निर्बंध} \mr{घालणे}, to impose restrictions.

\begin{cousinweb}
\textbf{\mr{निर्बंध}} (\emph{nirbandh}) is a tatsama, taken from Sanskrit \emph{nirbandha}, \textquotedblleft{}perseverance, insistence\textquotedblright{} (CDIAL 7355); Marathi uses it for a binding limit.
\end{cousinweb}

\begin{grammarlens}[title={What you've built}]
One more word for saying what happens if.
\end{grammarlens}

\subsection*{Your turn}
\begin{itemize}
\raggedright
  \item \emph{Say it:} \emph{nirbandh}
  \item \emph{Say it:} \emph{nirbandh}, once more
  \item \emph{Say it:} \emph{purāmuḷe pravāsāvar nirbandh ghātle āhet.} — Restrictions have been placed on travel because of the flood.
  \item \emph{Recall:} say the Marathi for a condition, then the Marathi for an exception, then say \emph{nirbandh} again
\end{itemize}

\begin{tcolorbox}[breakable,colback=teal!4,colframe=teal!35!black,title={Before you move on}]
What does \mr{निर्बंध} mean? (\textquotedblleft{}A curb; a restriction\textquotedblright{}.) Say it once more.
\end{tcolorbox}
\section[\texorpdfstring{\mr{निकष} (nikaṣ)}{निकष (nikaṣ)}]{\mr{निकष} (nikaṣ) — a criterion; a yardstick}

\label{lesson:MR-C378-nikas}



\begin{quote}
Before the new one: say the Marathi for an exception, then the Marathi for a curb.
\end{quote}

\subsection*{You'll want to know: \mr{निकष}}
\textbf{\mr{निकष}} — \emph{nikaṣ} — \textquotedblleft{}a criterion; a yardstick\textquotedblright{}.

\begin{quote}
\mr{शिष्यवृत्तीचे} \mr{निकष} \mr{या} \mr{वर्षी} \mr{बदलले}.
\end{quote}

\begin{quote}
\emph{śiṣyavṛttīce nikaṣ yā varṣī badalle.}
\end{quote}

\begin{quote}
The criteria for the scholarship changed this year.
\end{quote}

Masculine.

Plural \mr{निकष}; X \mr{चा} \mr{निकष}, the test for X.

\begin{cousinweb}
\textbf{\mr{निकष}} (\emph{nikaṣ}) is a tatsama, taken from Sanskrit \emph{nikaṣa}, \textquotedblleft{}roller; touchstone\textquotedblright{} (CDIAL 7139): the stone gold is tested on.
\end{cousinweb}

\begin{grammarlens}[title={What you've built}]
One more word for saying what happens if.
\end{grammarlens}

\subsection*{Your turn}
\begin{itemize}
\raggedright
  \item \emph{Say it:} \emph{nikaṣ}
  \item \emph{Say it:} \emph{nikaṣ}, once more
  \item \emph{Say it:} \emph{śiṣyavṛttīce nikaṣ yā varṣī badalle.} — The criteria for the scholarship changed this year.
  \item \emph{Recall:} say the Marathi for an exception, then the Marathi for a curb, then say \emph{nikaṣ} again
\end{itemize}

\begin{tcolorbox}[breakable,colback=teal!4,colframe=teal!35!black,title={Before you move on}]
What does \mr{निकष} mean? (\textquotedblleft{}A criterion; a yardstick\textquotedblright{}.) Say it once more.
\end{tcolorbox}
\section[\texorpdfstring{\mr{पूर्तता} \mr{करणे} (pūrtatā karṇe)}{पूर्तता करणे (pūrtatā karṇe)}]{\mr{पूर्तता} \mr{करणे} (pūrtatā karṇe) — to fulfil (a condition)}

\label{lesson:MR-C378-purtata-karne}



\begin{quote}
Before the new one: say the Marathi for a curb, then the Marathi for a criterion.
\end{quote}

\subsection*{You'll want to know: \mr{पूर्तता} \mr{करणे}}
\textbf{\mr{पूर्तता} \mr{करणे}} — \emph{pūrtatā karṇe} — \textquotedblleft{}to fulfil (a condition)\textquotedblright{}.

\begin{quote}
\mr{सगळ्या} \mr{अटींची} \mr{पूर्तता} \mr{केल्यास} \mr{प्रवेश} \mr{मिळेल}.
\end{quote}

\begin{quote}
\emph{sagaḷyā aṭīñcī pūrtatā kelyās praveś miḷel.}
\end{quote}

\begin{quote}
If all the conditions are fulfilled, you will get admission.
\end{quote}

Past, with -\mr{ने} (he …): \mr{त्याने} \mr{पूर्तता} \mr{केली}, \emph{tyāne pūrtatā kelī}.

What is fulfilled takes -\mr{ची}: \mr{अटींची} \mr{पूर्तता}.

\begin{cousinweb}
\textbf{\mr{पूर्तता} \mr{करणे}} (\emph{pūrtatā karṇe}) joins \mr{पूर्तता}, a learned noun built on Sanskrit \emph{pūrti}, \textquotedblleft{}filling, completion\textquotedblright{} (Monier-Williams), to \mr{करणे}, \textquotedblleft{}to do\textquotedblright{}.
\end{cousinweb}

\begin{grammarlens}[title={What you've built}]
One more word for saying what happens if.
\end{grammarlens}

\subsection*{Your turn}
\begin{itemize}
\raggedright
  \item \emph{Say it:} \emph{pūrtatā karṇe}
  \item \emph{Say it:} \emph{pūrtatā karṇe}, once more
  \item \emph{Say it:} \emph{sagaḷyā aṭīñcī pūrtatā kelyās praveś miḷel.} — If all the conditions are fulfilled, you will get admission.
  \item \emph{Recall:} say the Marathi for a curb, then the Marathi for a criterion, then say \emph{pūrtatā karṇe} again
\end{itemize}

\begin{tcolorbox}[breakable,colback=teal!4,colframe=teal!35!black,title={Before you move on}]
What does \mr{पूर्तता} \mr{करणे} mean? (\textquotedblleft{}To fulfil (a condition)\textquotedblright{}.) Say it once more.
\end{tcolorbox}
